\chapter{Introduction}
\label{ch:introduction}

% ─────────────────────────────────────────────────────────────────────────────
\section{Background and Motivation}
\label{sec:intro_background}
% ─────────────────────────────────────────────────────────────────────────────

The dominant paradigm in digital learning platforms remains one-size-fits-all
content delivery. A learner who already understands tectonic plate mechanics
receives the same sequence of questions as one encountering the concept for the
first time. This homogeneity is not a technical constraint but a design choice
that the empirical literature consistently identifies as suboptimal.

\textcite{bloom1984} demonstrated that one-to-one human tutoring, which adapts
to the learner's current ability in real time, produces learning outcomes
approximately two standard deviations above conventional classroom instruction.
\textcite{vanlehn2011} subsequently showed that step-based Intelligent Tutoring
Systems (ITS) can approach this performance at scale, achieving effect sizes
statistically indistinguishable from human tutoring. The adaptive feedback loop
is the critical variable, not the delivery medium.

Despite this evidence, adaptive ITS remain largely confined to mathematics and
programming domains, where knowledge structures are well-defined and assessment
is unambiguous. Geography, integrating declarative factual knowledge, procedural
reasoning, and conceptual understanding, has received almost no attention in the
adaptive learning literature \parencite{bednarz2013}. The UK KS3/GCSE Geography
curriculum represents both an underserved educational context and a meaningful
test case for ITS methods applied to non-STEM knowledge domains.

% ─────────────────────────────────────────────────────────────────────────────
\section{Research Question}
\label{sec:intro_rq}
% ─────────────────────────────────────────────────────────────────────────────

This project addresses the following research question:

\begin{quote}
\textit{Can a hybrid IRT 3PL + BKT adaptive engine, implemented in a
production web application, deliver personalised question sequencing that
produces measurable within-session learning gains for learners studying
UK KS3/GCSE Geography content?}
\end{quote}

The question has three components: an \textit{algorithmic claim} (the hybrid
IRT+BKT engine can perform valid adaptive selection), a \textit{technical
claim} (implementable in a production-deployed web application), and an
\textit{empirical claim} (the system produces measurable learning gains).
The evaluation in Chapter~\ref{ch:results} addresses all three, with the
acknowledgement that the empirical claim is evaluated at proof-of-concept
scale rather than at statistical significance.

% ─────────────────────────────────────────────────────────────────────────────
\section{Project Overview}
\label{sec:intro_overview}
% ─────────────────────────────────────────────────────────────────────────────

GeoMentor is an adaptive intelligent tutoring system for UK KS3/GCSE
Geography, accessible at \url{https://e-learning-final-project.vercel.app}.
The system implements a hybrid IRT 3PL + BKT algorithm that continuously
estimates learner ability and Knowledge Component mastery, selects questions
targeting the learner's Zone of Proximal Development, and enforces Bloom's
taxonomy cognitive prerequisites across a 13-KC curriculum graph.

The system was developed over ten two-week sprints from October 2025 to
May 2026, following a Design Science Research methodology \parencite{hevner2004}.
It is production-deployed, publicly accessible, and open-source, providing a
proof-of-concept implementation that can be independently verified and extended.

% ─────────────────────────────────────────────────────────────────────────────
\section{Contributions}
\label{sec:intro_contributions}
% ─────────────────────────────────────────────────────────────────────────────

This project makes the following contributions:

\begin{enumerate}[leftmargin=2em]

  \item \textbf{A production-deployed hybrid IRT 3PL + BKT adaptive engine
  for UK Geography.} The first publicly accessible ITS targeting the UK
  KS3/GCSE Geography curriculum, extending the empirical evidence base for
  adaptive learning beyond traditional STEM domains.

  \item \textbf{A credit-factor modification to the standard BKT update.}
  A novel modification of the BKT learning transition equation that weights
  correct responses by a graduated credit factor reflecting the maximum hint
  level accessed, formalising the evidential relationship between scaffolding
  usage and mastery evidence.

  \item \textbf{Bloom-gated adaptive progression.} An implementation of
  Bloom's taxonomy as a computational selection constraint, enforcing that
  learners consolidate recall-level knowledge before encountering
  application-level items.

  \item \textbf{A five-track multi-method evaluation framework for small-sample
  ITS validation.} A structured evaluation approach combining algorithmic
  validation, engagement data, usability metrics, misconception analysis, and
  robustness monitoring, designed for contexts where inferential statistics are
  not viable due to small sample sizes.

  \item \textbf{An open-source domain-agnostic ITS architecture.} The codebase
  provides a reusable template for researchers and developers building adaptive
  systems for non-STEM declarative-knowledge domains.

\end{enumerate}

% ─────────────────────────────────────────────────────────────────────────────
\section{Scope and Delimitations}
\label{sec:intro_scope}
% ─────────────────────────────────────────────────────────────────────────────

\begin{itemize}[leftmargin=2em]
  \item \textbf{Domain.} UK KS3/GCSE Geography only (AQA specification).
  \item \textbf{Assessment type.} Multiple-choice questions with four options.
  \item \textbf{Participant population.} University of Hull students
    (convenience sample); the system targets KS3/GCSE students but was
    evaluated with adults due to ethical and logistical constraints.
  \item \textbf{Algorithm.} IRT 3PL + BKT hybrid. DKT, Elo, and 2PL were
    evaluated and rejected (Section~\ref{sec:lr_algorithms}).
  \item \textbf{Evaluation scale.} Proof-of-concept at $n = 4$. No
    inferential statistical claims are made.
\end{itemize}

% ─────────────────────────────────────────────────────────────────────────────
\section{Ethical and Professional Considerations}
\label{sec:intro_ethics}
% ─────────────────────────────────────────────────────────────────────────────

All participant data was collected under informed written consent. No
personally identifiable data beyond email address and name (required for
authentication) was collected; session performance data is stored under
a pseudonymous user ID, accessible only to the authenticated user and
to the researcher under the study protocol.

\textbf{AI use declaration.} This project was developed with AI assistance
(Claude, Anthropic), consistent with the University of Hull's Student AI Use
Policy (2025). AI assistance was employed for infrastructure scaffolding
(Supabase RLS boilerplate, Resend email integration, bcrypt helpers), frontend
component implementation from candidate-authored specifications, Mapbox
integration, Vercel deployment configuration, and Prisma schema iteration.
The following constitute the candidate's independent intellectual work and
were not produced by AI tools: the hybrid IRT 3PL + BKT algorithm architecture;
the credit-factor BKT modification and its full mathematical derivation; the
Bloom escalation gating design and its implementation as a computational
constraint; the 13-KC curriculum ontology and prerequisite DAG; the research
design including the within-subjects quasi-experimental structure, VanLehn
benchmark selection, and Sprint~5 scope reduction justification; the five-track
multi-method evaluation framework; all Architecture Decision Records; and the
IRT calibration strategy. AI assistance was employed exclusively as an
implementation accelerator for repetitive or infrastructural tasks; all design
decisions and algorithmic contributions are the candidate's own work and are
grounded in the peer-reviewed literature cited throughout this report.

Data collection complied with UK GDPR and the Data Protection Act 2018.
The lawful basis for processing is \textit{consent}: participants provided
explicit informed written consent prior to any data collection and were
informed of their unconditional right to withdraw without consequence.
No special category data were collected; all data are stored on EU-resident
servers (Supabase, Frankfurt region). This project was conducted in accordance
with the British Computer Society Code of Conduct's four obligations (public
interest, professional competence, integrity, and duty to relevant authority);
ethical considerations were assessed under the University of Hull Research
Ethics Policy prior to participant recruitment.

% ─────────────────────────────────────────────────────────────────────────────
\section{Report Structure}
\label{sec:intro_structure}
% ─────────────────────────────────────────────────────────────────────────────

Chapter~\ref{ch:literature_review} critically reviews adaptive ITS theory,
algorithms, CAT principles, empirical effectiveness, and ethical dimensions,
identifying five research gaps. Chapter~\ref{ch:methodology} details the
KC ontology, adaptive algorithm design with full mathematical derivation,
evaluation protocol, and agile development methodology.
Chapter~\ref{ch:implementation} describes the technical implementation
including database schema, adaptive engine, API routes, and security.
Chapter~\ref{ch:results} presents the five-track multi-method evaluation.
Chapter~\ref{ch:discussion} interprets results, examines limitations, and
identifies future work. Chapter~\ref{ch:conclusion} summarises contributions
and offers a final statement.
